% !TEX root = ../main.tex
%% Corrigé du sujet type bac n° 8 — résultats vérifiés par programme (bac_s8.py)

\corrigeSujet{8}

\begin{corrigeBac}{4}{Questionnaire à choix multiples}
\begin{questions}
  \item \textbf{Réponse b).} Il y a $5!=120$ permutations des cinq lettres, mais échanger les deux E donne le même mot :
    $\frac{120}{2}=60$ anagrammes. \pts{1 pt}
  \item \textbf{Réponse c).} $P(A\cap B)=0{,}4\times0{,}5=0{,}2$, d'où $P(A\cup B)=0{,}4+0{,}5-0{,}2=0{,}7$. \pts{1 pt}
  \item \textbf{Réponse d).} $np=12$ et $np(1-p)=3$ donnent $1-p=\frac14$, soit $p=0{,}75$, puis $n=16$. \pts{1 pt}
  \item \textbf{Réponse b).} L'ordre ne compte pas : $\binom{10}{3}=120$. \pts{1 pt}
\end{questions}
\end{corrigeBac}

\begin{corrigeBac}{5}{Suites et récurrence}
\begin{questions}
  \item $u_1=\frac12+0+1=\res{\tfrac32}$, $u_2=\frac34+1+1=\res{\tfrac{11}4}$ et $u_3=\frac{11}8+2+1=\res{\tfrac{35}{8}}$. \pts{0,5 pt}
  \item \etape{Initialisation} $u_0=1\geq0$.\\
    \etape{Hérédité} Si $u_n\geq n$, alors $u_{n+1}=\frac12u_n+n+1\geq\frac12n+n+1\geq n+1$.\\
    \etape{Conclusion} Pour tout entier naturel $n$, $u_n\geq n$. Comme $\lim n=+\infty$, le théorème de comparaison donne
    $\res{\lim u_n=+\infty}$. \pts{1,5 pt}
  \item $v_{n+1}=u_{n+1}-2(n+1)+2=\frac12u_n+n+1-2n=\frac12u_n-n+1=\frac12\left(u_n-2n+2\right)=\frac12v_n$ : la suite $(v_n)$ est
    géométrique de raison $\res{\tfrac12}$, de premier terme $v_0=1-0+2=\res{3}$. \pts{1,25 pt}
  \item $v_n=3\times\left(\frac12\right)^{n}$, d'où $u_n=v_n+2n-2=3\times\left(\frac12\right)^{n}+2n-2$. \pts{0,5 pt}
  \item Pour $n\geq1$, $\dfrac{u_n}{n}=\dfrac{3\times0{,}5^{n}}{n}+2-\dfrac2n$. Les deux fractions tendent vers $0$ :
    $\res{\lim\dfrac{u_n}{n}=2}$. \pts{0,75 pt}
  \item On accumule les termes en les calculant de proche en proche.
\begin{codecote}[0.45]
\begin{pythoncol}
def somme(n):
    u = 1
    S = u
    for k in range(n):
        u = 0.5 * u + k + 1
        S = S + u
    return S
\end{pythoncol}
\cote
\code{somme(10)} renvoie environ \code{93.997} : $\res{S_{10}\approx93{,}997}$ (la valeur exacte est $\frac{96\,253}{1\,024}$).
\end{codecote}
    \pts{0,5 pt}
\end{questions}
\end{corrigeBac}

\begin{corrigeBac}{6}{Fonction exponentielle et convexité}
\partieExo{A}{Étude de la fonction}
\begin{questions}
  \item En $-\infty$ : $x^{2}e^{x}\to0$ (croissances comparées) et $e^{x}\to0$, d'où $\res{\lim\limits_{x\to-\infty}f(x)=0}$ : l'axe des
    abscisses est asymptote en $-\infty$. En $+\infty$ : $x^{2}-3\to+\infty$ et $e^{x}\to+\infty$, d'où $\res{\lim\limits_{x\to+\infty}f(x)=+\infty}$. \pts{1 pt}
  \item $f'(x)=2x\,e^{x}+\left(x^{2}-3\right)e^{x}=\left(x^{2}+2x-3\right)e^{x}=(x-1)(x+3)e^{x}$, du signe de $(x-1)(x+3)$.
\begin{center}
\begin{tableau}
  \otTab[1.8]{$x$/0.8,$f'(x)$/0.8,$f$/1.6}{$-\infty$,$-3$,$1$,$+\infty$}
  \tkzTabLine{,+,z,-,z,+,}
  \tkzTabVar{-/$0$,+/$6e^{-3}$,-/$-2e$,+/$+\infty$}
\end{tableau}
\end{center}
    Maximum local $f(-3)=6e^{-3}\approx0{,}30$ ; minimum $f(1)=-2e\approx-5{,}44$. \pts{1,5 pt}
  \item $f(x)=0\iff x^{2}=3\iff\res{x=-\sqrt3\text{ ou }x=\sqrt3}$ (car $e^{x}>0$). $f$ est positive sur $\intOF{-\infty}{-\sqrt3}$ et sur
    $\intFO{\sqrt3}{+\infty}$, négative sur $\intFF{-\sqrt3}{\sqrt3}$. \pts{0,5 pt}
\end{questions}
\partieExo{B}{Convexité et aire}
\begin{questions}[resume]
  \item $f''(x)$ est du signe de $x^{2}+4x-1$, de racines $-2\pm\sqrt5$ : $f$ est convexe sur $\intOF{-\infty}{-2-\sqrt5}$ et sur
    $\intFO{-2+\sqrt5}{+\infty}$, concave entre les deux. Les points d'inflexion ont pour abscisses $\res{-2-\sqrt5}$ et $\res{-2+\sqrt5}$. \pts{1 pt}
  \item $F'(x)=(2x-2)e^{x}+\left(x^{2}-2x-1\right)e^{x}=\left(x^{2}-3\right)e^{x}=f(x)$. \pts{0,75 pt}
  \item Sur $\intFF{-\sqrt3}{\sqrt3}$, $f\leq0$ : l'aire vaut $\mathcal{A}=-\int_{-\sqrt3}^{\sqrt3}f(x)\,\mathrm{d}x=F\left(-\sqrt3\right)-F\left(\sqrt3\right)$.
    Or $F\left(\sqrt3\right)=\left(2-2\sqrt3\right)e^{\sqrt3}$ et $F\left(-\sqrt3\right)=\left(2+2\sqrt3\right)e^{-\sqrt3}$, d'où
    $\mathcal{A}=\res{\left(2+2\sqrt3\right)e^{-\sqrt3}+\left(2\sqrt3-2\right)e^{\sqrt3}}\approx\res{9{,}24}$ u.a. \pts{1,25 pt}
\end{questions}
\end{corrigeBac}

\begin{corrigeBac}{5}{Géométrie : distance à une droite}
\begin{questions}
  \item $\res{d:\ (1+t\pv t\pv2-t)}$, $t\in\R$. \pts{0,5 pt}
  \item $\vect{MM_t}\,(t-2\pv t-4\pv1-t)$, donc $MM_t^{2}=(t-2)^{2}+(t-4)^{2}+(1-t)^{2}=3t^{2}-14t+21$. \pts{1 pt}
  \item Le trinôme $3t^{2}-14t+21$ est minimal pour $t=\frac{14}{6}=\frac73$ : $\res{H\left(\tfrac{10}{3}\pv\tfrac73\pv-\tfrac13\right)}$. La valeur
    minimale de $MM_t^{2}$ est $3\times\frac{49}9-14\times\frac73+21=\frac{14}{3}$, d'où la distance minimale $\res{\sqrt{\tfrac{14}{3}}=\tfrac{\sqrt{42}}{3}}\approx2{,}16$. \pts{1,25 pt}
  \item $\vect{MH}\left(\frac13\pv-\frac53\pv-\frac43\right)$ et $\vect{MH}\cdot\vec u=\frac13-\frac53+\frac43=0$ : $(MH)\perp d$. \pts{0,5 pt}
  \item $\vec u$ est normal à $\mathcal{Q}$ : $x+y-z+d=0$, et $M$ donne $d=-6$ : $\res{\mathcal{Q}:\ x+y-z-6=0}$. Avec la représentation
    de $d$ : $(1+t)+t-(2-t)-6=3t-7=0$, d'où $t=\frac73$ : c'est le point $H$. \pts{1 pt}
  \item $MM_t=\sqrt6\iff3t^{2}-14t+21=6\iff3t^{2}-14t+15=0$, de discriminant $16$ : $t=3$ ou $t=\frac53$. Les points cherchés sont
    $\res{(4\pv3\pv-1)}$ et $\res{\left(\tfrac83\pv\tfrac53\pv\tfrac13\right)}$. \pts{0,75 pt}
\end{questions}
\end{corrigeBac}
